\begin{table}[H]
\centering
\caption{Nama Daerah Perbandingan Cluster Biplot vs Kategori IPM (UNDP) (2022)} \label{tab:app2crossname2022}
\begin{adjustbox}{max width=\textwidth}
\begin{tabular}{|p{4cm}|p{10cm}|}
\hline
\textbf{Kategori} & \textbf{Daftar Wilayah} \\
\hline
Cl. 1 (Rendah) \& Sangat Tinggi & - \\
\hline
Cl. 1 (Rendah) \& Tinggi & Kep. Seribu, Ciamis, Kota Cilegon \\
\hline
Cl. 1 (Rendah) \& Sedang & Sukabumi, Cianjur, Garut, Tasikmalaya, Majalengka, Indramayu, Pangandaran, Pandeglang, Lebak, Serang \\
\hline
Cl. 1 (Rendah) \& Rendah & - \\
\hline
Cl. 2 (Sedang) \& Sangat Tinggi & - \\
\hline
Cl. 2 (Sedang) \& Tinggi & Bogor, Bandung, Kuningan, Cirebon, Sumedang, Purwakarta, Karawang, Kota Sukabumi, Kota Cirebon, Kota Tasikmalaya, Kota Banjar, Tangerang, Kota Serang \\
\hline
Cl. 2 (Sedang) \& Sedang & Subang, Bandung Barat \\
\hline
Cl. 2 (Sedang) \& Rendah & - \\
\hline
Cl. 3 (Tinggi) \& Sangat Tinggi & Kota Jakarta Selatan, Kota Jakarta Timur, Kota Jakarta Pusat, Kota Jakarta Barat, Kota Jakarta Utara, Kota Bandung, Kota Bekasi, Kota Depok, Kota Tangerang Selatan \\
\hline
Cl. 3 (Tinggi) \& Tinggi & Bekasi, Kota Bogor, Kota Cimahi, Kota Tangerang \\
\hline
Cl. 3 (Tinggi) \& Sedang & - \\
\hline
Cl. 3 (Tinggi) \& Rendah & - \\
\hline
\end{tabular}
\end{adjustbox}
\end{table}